
\medskip
\noindent\textbf{4.3.}
\phantomsection
\label{sga3:VIIA:4.3}

Finally, let \(G\) be an \(S\)-group scheme in abelian groups. The
composite morphism
\[
  \mu(G):
  U^p(G)\lhook\joinrel\longrightarrow G_S^p
  \longrightarrow G,
\]
defined by multiplication, factors through \(V^p(G)\). Thus there
is a morphism
\[
  \nu(G):V^p(G)\longrightarrow G
\]
such that
\[
  \nu(G)\circ q'(G)=\mu(G).
\]
The resulting commutative diagram is
\begin{center}
  \phantomsection
  \label{sga3:VIIA:4.3-verschiebung-master-diagram}
  \begin{tikzpicture}[
  baseline=(current bounding box.center),
  >=Stealth,
  every node/.style={font=\normalsize}
]
  \node (gleft) at (-3.3,2.1) {$G$};
  \node (upg)   at (0,2.1) {$\mathrm U^p(G)$};
  \node (gright) at (3.3,2.1) {$G$};
  \node (vpg)   at (-1.65,0.35) {$\mathrm V^p(G)$};
  \node (gbr)   at (1.65,0.35) {$G^{[p/S]}$};
  \node (gps)   at (0,-1.65) {$G^{(p/S)}$};

  \draw[->] (upg) -- node[above] {$\mu(G)$} (gleft);
  \draw[->] (gright) -- node[above] {$\delta(G)$} (upg);
  \draw[->] (vpg) -- node[pos=.52,above,fill=white,inner sep=1pt]
    {$\nu(G)$} (gleft);
  \draw[->] (upg) -- node[pos=.53,right,fill=white,inner sep=1pt]
    {$q'(G)$} (vpg);
  \draw[->] (gbr) -- node[below] {$i(G)$} (vpg);
  \draw[->] (gright) -- node[pos=.5,left,fill=white,inner sep=1pt]
    {$F^{[p]}(G/S)$} (gbr);
  \draw[->] (gbr) -- node[pos=.55,anchor=west,xshift=2pt,fill=white,inner sep=1pt]
    {$\varphi(G)$} (gps);
  \draw[->] (gright) to[out=-92,in=18]
    node[right] {$\operatorname{Fr}(G/S)$} (gps);
  \draw[->,dotted] (gps) to[out=172,in=-87]
    node[left] {$\operatorname{Ver}(G/S)$} (gleft);
  \node at (3.55,-1.65) {$.$};
\end{tikzpicture}

\end{center}

When \(G\) is flat over \(S\), \(\varphi(G)\) is an isomorphism,
and we may define the morphism called the \emph{Verschiebung}
\[
  \operatorname{Ver}(G/S):
  G^{(p/S)}
  \longrightarrow
  G
\]
by
\[
  \operatorname{Ver}(G/S)
  =
  \nu(G)\circ i(G)\circ\varphi(G)^{-1}.
\]
As \(G\) ranges over the flat commutative \(S\)-group schemes, the
assignment
\[
  \operatorname{Ver}:G\longmapsto\operatorname{Ver}(G/S)
\]
is plainly a functorial homomorphism. Consequently,
\(\operatorname{Ver}(G/S)\) is a group homomorphism.

For every \(S\)-scheme \(T\), the composite
\[
  G(T)
  \xrightarrow{\delta(G)(T)}
  U^p(G)(T)
  \xrightarrow{\mu(G)(T)}
  G(T)
\]
sends \(x\in G(T)\) to \(p\cdot x\). We may therefore write
\(p\cdot\operatorname{id}_G\) in place of
\(\mu(G)\circ\delta(G)\), obtaining the classical formula
\[
  \tag{*}\label{sga3:vii:paragraph:4-3:equation:star}
  \operatorname{Ver}(G/S)\circ\operatorname{Fr}(G/S)
  =
  p\cdot\operatorname{id}_G.
\]

\medskip
\noindent\textbf{Examples 4.3.1.}
\phantomsection
\label{sga3:VIIA:4.3.1}

\smallskip
\noindent\textit{a)}
If \(G\) is a constant \(S\)-scheme in abelian groups, then
\(\operatorname{Fr}(G/S)\) identifies with the identity morphism
of \(G\) (cf.\ \SGARefLink{sga3:VIIA:4.1.1.1}{4.1.1.1}). Hence
\[
  \operatorname{Ver}(G/S)
  =
  p\cdot\operatorname{id}_G.
\]

\smallskip
\noindent\textit{b)}
If \(G\) is the diagonalizable \(S\)-group of type \(M\), then
\(\operatorname{Fr}(G/S)=p\cdot\operatorname{id}_G\) by
\SGARefLink{sga3:VIIA:4.1.4}{4.1.4}~\textit{a)}. It follows at once that
\(\operatorname{Ver}(G/S)\) is the identity morphism of \(G\).

\smallskip
\noindent\textit{c)}
If \(\mathcal E\) is a flat \(\mathcal O_S\)-module and
\(G=\mathcal V(\mathcal E)\), then
\(\operatorname{Ver}(G/S)\) is zero, as is
\(p\cdot\operatorname{id}_G\). In Expos\'e VII\textsubscript{B} we
shall see that a commutative algebraic group \(G\) over a field
\(k\) is ``unipotent'' if and only if, for some \(n\), the
composite homomorphism
\begin{center}
  \phantomsection
  \label{sga3:VIIA:4.3.1-iterated-verschiebung-chain}
  \(\displaystyle
G^{(p^n)}
  \xrightarrow{\operatorname{Ver}(G^{(p^{n-1})}/S)}
G^{(p^{n-1})}
  \longrightarrow \cdots \longrightarrow
G^{(p)}
  \xrightarrow{\operatorname{Ver}(G/S)}
G
\)

\end{center}
is zero. Here
\[
  G^{(p^n)}
  =
  \bigl(G^{(p^{n-1})}\bigr)^{(p)}
\]
as in \SGARefLink{sga3:VIIA:4.1.3}{4.1.3}.

\medskip
\noindent\textbf{4.3.2.}
\phantomsection
\label{sga3:VIIA:4.3.2}

Since
\(\operatorname{Ver}:G\mapsto\operatorname{Ver}(G/S)\) is a
functorial homomorphism on the category of flat commutative
\(S\)-group schemes, the square
\begin{center}
  \phantomsection
  \label{sga3:VIIA:4.3.2-frobenius-verschiebung-square}
  \begin{tikzcd}[row sep=3.3em,column sep=5.2em]
G^{(p)} \arrow[r,"\operatorname{Ver}(G/S)"]
  \arrow[d,swap,"\operatorname{Fr}(G/S)^{(p)}"]
  & G \arrow[d,"\operatorname{Fr}(G/S)"] \\
G^{(p^2)} \arrow[r,"\operatorname{Ver}(G^{(p)}/S)"]
  & G^{(p)}
\end{tikzcd}

\end{center}
commutes. Here \(\operatorname{Fr}(G/S)^{(p)}\) denotes the
pullback of \(\operatorname{Fr}(G/S)\) along the base change
\(\operatorname{fr}(S)\). By \SGARefLink{sga3:VIIA:4.1.1}{4.1.1},
\[
  \operatorname{Fr}(G/S)^{(p)}
  =
  \operatorname{Fr}(G^{(p)}/S).
\]
Applying formula \((*)\) of \SGARefLink{sga3:VIIA:4.3}{4.3} to \(G^{(p)}\), we obtain
\[
  \tag{**}\label{sga3:vii:paragraph:4-3-2:equation:double-star}
  \operatorname{Fr}(G/S)\circ\operatorname{Ver}(G/S)
  =
  \operatorname{Ver}(G^{(p)}/S)
    \circ\operatorname{Fr}(G^{(p)}/S)
  =
  p\cdot\operatorname{id}_{G^{(p)}}.
\]

\medskip
\noindent\textbf{4.3.3.}
\phantomsection
\label{sga3:VIIA:4.3.3}

Suppose finally that \(G\) is a commutative \(S\)-group which is
finite and locally free. Let \(\mathcal A\) be its affine
\(\mathcal O_S\)-algebra, and let \(\pi\) be the endomorphism of
the sheaf of rings \(\mathcal O_S\) which sends a section \(x\) to
\(x^p\).\footnote{The editors changed the order of presentation,
introducing first the objects which occur in the diagram below.}
Denote by \(\Sigma^p\mathcal A\) the subalgebra of
\(\bigotimes_{\mathcal O_S}^p\mathcal A\) consisting of the
sections invariant under the symmetric-group action, and by
\[
  i(\mathcal A):
  \Sigma^p\mathcal A
  \longrightarrow
  \bigotimes_{\mathcal O_S}^p\mathcal A
\]
its inclusion.

Let
\[
  \Delta^p(\mathcal A):
  \mathcal A
  \longrightarrow
  \bigotimes_{\mathcal O_S}^p\mathcal A
\]
be the morphism obtained by iterating the diagonal morphism of the
coalgebra \(\mathcal A\). It corresponds to the multiplication
morphism \(U^p(G)=G_S^p\to G\). By the beginning of \SGARefLink{sga3:VIIA:4.3}{4.3},
\(\Delta^p(\mathcal A)\) factors through
\(\Sigma^p\mathcal A\); that is, it induces a morphism
\[
  a(\mathcal A):
  \mathcal A\longrightarrow\Sigma^p\mathcal A
\]
such that
\[
  i(\mathcal A)\circ a(\mathcal A)
  =
  \Delta^p(\mathcal A).
\]

On the other hand, let \(S^p(\mathcal A)\) be the homogeneous
component of degree \(p\) in the symmetric algebra of
\(\mathcal A\), and let
\[
  q(\mathcal A):
  \bigotimes_{\mathcal O_S}^p\mathcal A
  \longrightarrow
  S^p(\mathcal A)
\]
be the canonical projection. The multiplication
\[
  m^p(\mathcal A):
  \bigotimes_{\mathcal O_S}^p\mathcal A
  \longrightarrow
  \mathcal A
\]
factors through \(S^p(\mathcal A)\); equivalently, it induces a map
\[
  b(\mathcal A):
  S^p(\mathcal A)\longrightarrow\mathcal A
\]
such that
\[
  b(\mathcal A)\circ q(\mathcal A)
  =
  m^p(\mathcal A).
\]

Since \(\Sigma^p\mathcal A\) is the affine algebra of
\(V^p(G)\), the composite
\(i(G)\circ\varphi(G)^{-1}\) induces, again by the beginning of
\SGARefLink{sga3:VIIA:4.3}{4.3}, an algebra homomorphism
\[
  r(\mathcal A):
  \Sigma^p\mathcal A
  \longrightarrow
  \mathcal A\otimes_\pi\mathcal O_S.
\]
It vanishes on sections of the form
\[
  \sum_{\sigma\in\mathfrak S_p}
    a_{\sigma(1)}\otimes\cdots\otimes a_{\sigma(p)}
\]
and sends \(a\otimes\cdots\otimes a\) to
\(a\otimes_\pi1\). Similarly, \(j(\mathcal A)\) is the morphism
of \(\mathcal O_S\)-modules
\[
  a\otimes_\pi1
  \longmapsto
  q(a\otimes\cdots\otimes a).
\]

These maps form the commutative diagram
\begin{center}
  \phantomsection
  \label{sga3:VIIA:4.3.3-cartier-duality-algebra-diagram}
  \begin{tikzpicture}[
  baseline=(current bounding box.center),
  >=Stealth,
  every node/.style={font=\normalsize}
]
  \node (atl) at (-3.2,2.0) {$\mathcal A$};
  \node (atm) at (0,2.0) {$\displaystyle\bigotimes_{\mathcal O_S}^{p}\mathcal A$};
  \node (atr) at (3.2,2.0) {$\mathcal A$};
  \node (aml) at (-1.65,0.15) {$\Sigma^p\mathcal A$};
  \node (amr) at (1.65,0.15) {$\mathrm S^p(\mathcal A)$};
  \node (ab)  at (0,-1.75) {$\mathcal A\otimes_{\pi}\mathcal O_S$};
  \node at (-4.85,0.15) {$({\mathcal A})$};

  \draw[->] (atl) -- node[above] {$\Delta^p(\mathcal A)$} (atm);
  \draw[->] (atm) -- node[above] {$m^p(\mathcal A)$} (atr);
  \draw[->] (atl) -- node[midway,left,fill=white,inner sep=1pt]
    {$a(\mathcal A)$} (aml);
  \draw[->] (aml) -- node[midway,left,fill=white,inner sep=1pt]
    {$i(\mathcal A)$} (atm);
  \draw[->] (atm) -- node[midway,right,fill=white,inner sep=1pt]
    {$q(\mathcal A)$} (amr);
  \draw[->] (amr) -- node[midway,right,fill=white,inner sep=1pt]
    {$b(\mathcal A)$} (atr);
  \draw[->] (aml) -- node[midway,left,fill=white,inner sep=1pt]
    {$r(\mathcal A)$} (ab);
  \draw[->] (ab) -- node[midway,right,fill=white,inner sep=1pt]
    {$j(\mathcal A)$} (amr);
  \node at (3.55,-1.75) {$.$};
\end{tikzpicture}

\end{center}
The composite \(r(\mathcal A)\circ a(\mathcal A)\) corresponds to
the Verschiebung morphism \(\operatorname{Ver}(G/S)\), while
\(b(\mathcal A)\circ j(\mathcal A)\) corresponds to the Frobenius
morphism \(\operatorname{Fr}(G/S)\).

The preceding commutative diagram \((\mathcal A)\) is self-dual.
Indeed, let \(D\) be the functor which associates to each
\(\mathcal O_S\)-module \(\mathcal M\) its dual
\[
  D\mathcal M
  =
  \operatorname{Hom}_{\mathcal O_S}
    (\mathcal M,\mathcal O_S).
\]
The image of diagram \((\mathcal A)\) under \(D\) is precisely
diagram \((D\mathcal A)\). The morphisms
\[
  Dr(\mathcal A),\quad
  Da(\mathcal A),\quad
  Dj(\mathcal A),\quad
  Db(\mathcal A)
\]
identify respectively with
\[
  j(D\mathcal A),\quad
  b(D\mathcal A),\quad
  r(D\mathcal A),\quad
  a(D\mathcal A).
\]
By \SGARefLink{sga3:VIIA:3.3.1}{3.3.1}, it follows that:

\begin{quote}
\textit{In the category of commutative \(S\)-groups which are
finite and locally free, Cartier duality interchanges the
Frobenius morphism and the Verschiebung.}
\end{quote}

